\subsection{Approaches with increasing expressiveness}
\label{subsection_comparison_of_approaches}

\subsubsection{Classic approach vector algebra ($V\!A$)}

\begin{itemize}
    \item Use of scalars and vectors, exclusively
    \item Use of the cross-product between vectors to create directions perpendicular to
    both vectors (a concept that can only be defined in three-dimensional space)
    \item Rotations are defined by a rotation axis (a concept that is only defined in
    three-dimensional space)\footnote{The last two concepts cannot be extended to
    higher-dimensional spaces.}
\end{itemize}

\subsubsection{General observations for Geometric Algebra ($GA$)}

\begin{itemize}
    \item The GA approach is coordinate-free!
    \item W.r.t. a vector, the dot product (or interior product) takes into account the
    parts parallel to the vector, exclusively. Thus, it is not invertible itself.
    \item W.r.t. a vector, the wedge product (or outer product) takes into account the
    parts perpendicular to the vector, exclusively. Thus, it is not invertible itself.
    \item Only the geometric product taking both parts into account is invertible.
\end{itemize}

\subsubsection{Euclidean geometric algebra ($EGA$)}

\begin{itemize}
    \item Use of scalars, vectors, bivectors, ..., pseudoscalars (higher grade objects up
    to order $n$ depend on the dimension $n$ of the modelled space).
    \item Use of the outer product (wedge product) between vectors, which defines the
    space spanned by both vectors (can be defined in any dimension, applied as join
    operation).
    \item Use of the regressive outer product (regressive wedge product) which defines the
    space shared by both objects (represents the intersection of the spaces used by the
    objects, applied as meet operation).
    \item Modelled objects cover scalars, directions, planes (in $3D$) and the whole space
    directly.
    \item All objects modelled in the algebra contain the (non-distinguishable) origin
    $O=(0,0,...,0)$. Objects outside the origin must be modeled by explicitly adding
    translations (leading to objects that cannot directly be manipulated as easily as
    objects that directly represent subspaces within the algebra).
    \item Orthogonal transformations in the algebra can directly model reflections and
    rotations. Rotations are defined in a plane spanned by two vectors defining the
    rotation.
\end{itemize}

\subsubsection{Projective (Euclidean) geometric algebra ($PGA$)}

\begin{itemize}
    \item Like $EGA$, but in addition having a distinguishable origin encoded in an
    additional dimension $O=(0,0,...,0,1)$ (in a projective space with an additional
    dimension projected onto a fixed value, typically 1).
    \item Use of the outer and regressive outer products as join and meet operations.
    \item Directly models points, directions and planes (in $3D$), as well as the whole
    space, including objects located in the horizon, i.e. at infinity (points at
    infinity are modeled as directions towards those points).
    \item Can directly model higher grade objects by products between objects in the
    algebra that do not contain the origin, e.g. the wedge product between two arbitrary
    points defines a line.
    \item Orthogonal transformations in the algebra can directly model reflections,
    rotations and translations.
\end{itemize}

\subsubsection{Space-time algebra ($STA$)}

\begin{itemize}
    \item Extends the Euclidean approach to space-time: $\mathbb{G}(1,3,0)$ with three
    space-like and one time-like basis vector (Minkowski signature).
    \item Bivectors carry both kinds of transformation generators: space-like planes
    generate rotations, time-like planes generate Lorentz boosts. The even-grade rotors
    implement both through the same sandwich product as in $EGA$.
    \item Quantities that are separate objects in $3D$ combine into single space-time
    objects (events and four-velocities as vectors; the electric and magnetic fields as
    one field bivector), and frame changes act on all of them by rotor conjugation.
    \item Suitable when time is part of the modelled physics itself: fields of moving
    sources, induction, waves, invariance arguments. For static and quasi-static
    problems $EGA$ with time as an external parameter remains sufficient.
\end{itemize}

\subsubsection{Choosing the algebra}

The algebras are layers with distinct responsibilities rather than competing
formalisms. $EGA$ describes the \emph{value} of a physical quantity at a point: its
metric is invertible, which provides norms, angles and gradients, and with them field
calculus. Its objects carry no position -- the evaluation point is a parameter, not
part of the object. $PGA$ describes \emph{located} objects -- points, lines of action,
planes, poses -- and rigid motion; its degenerate metric intentionally gives up field
calculus to gain translations as sandwich products. $STA$ describes quantities whose
physics couples space and time. \\

A quantity therefore belongs in $EGA$ when its position is merely where it is evaluated
(a field value, an energy density), and in $PGA$ when its position is part of what it
is (a force on its line of action, the pose of a body). The two connect through the
ideal sector: an $EGA$ direction embeds as an ideal point, an $EGA$ bivector as an
ideal line (a couple), so a local result computed in $EGA$ is bound to its location by
a single wedge product -- e.g. $P \wedge f$ turns a force vector evaluated at a point
into the wrench acting on that point's line of action. $STA$ enters when time can no
longer stay an external parameter. \\

In short: $EGA$ answers \emph{what} -- the value of a quantity at a point. $PGA$
answers \emph{where} -- location and rigid motion. $STA$ answers \emph{when} --
coupling across space and time.

\newpage